This version of \textit{The Trouble at Round Rock}, told by Howard Gorman, provides a different perspective on the historical clash involving Black Horse. Gorman begins by recalling the cultural tradition of elders instructing youth and frames the story as both a historical account and a moral lesson. In this version, Gorman tells the story as reported by his grandfather (``Man Who Lassoes"), who calls Black Horse ``Butcher Squeezed Together In the Middle". Here, Black Horse is depicted as more boastful and antagonistic than in Left-Handed Mexican Clansman's telling. Gorman's grandfather confronts Black Horse directly and condemns his reckless actions and discouraging unnecessary conflict.

Unlike the earlier version, which portrays Black Horse's actions in a more collaborative resistance against the Agent, Gorman's account casts Black Horse as a figure of bravado who retreats when challenged. The narrative also shows how futile it is to cause disturbances, as other leaders are mentioned who gained nothing but hardship from similar conflicts. Through his grandfather's words \textit{T’áá t’óó ’ák’ijį’ dahojooł’chįįd. Jó ’akon, doo nijichxǫ’ da she’awéé’} (`They merely brought hardship upon themselves. So, don't sulk, my baby.'), Gorman conveys a moral message whereby restraint and peace should be favored over impulsive rebellion.
